% chapters/11/contoh.tex
\section{Dua Contoh Analisis}
\label{sec:contoh-bab-11}

\subsection{Contoh Pertama: Menafsirkan Tabel Evaluasi}

\begin{contohbox}
Pada Tabel~\ref{tab:evaluasi-selang}, pencarian setiap bingkai menjalankan
2.700 pencarian dengan jarak rata-rata 35 piksel, sedangkan selang 0,30 detik
menjalankan 858 pencarian dengan jarak rata-rata 34 piksel.
\end{contohbox}

Perbandingan itu dapat dinyatakan dalam bentuk rasio, sehingga lebih mudah
ditafsirkan.

\[
\text{rasio beban} = \frac{2700}{858} \approx 3{,}15,
\qquad
\text{selisih mutu} = 35 - 34 = 1 \text{ piksel}.
\]

Pekerjaan perhitungannya menjadi lebih dari tiga kali lipat, sedangkan
perbaikan mutunya hanya satu piksel, yaitu sekitar tiga persen dari jaraknya.
Kesimpulannya jelas: \textbf{pekerjaan tambahan itu tidak sepadan}.

Sebaliknya, membandingkan baris terakhir dengan baris ketiga memberikan hasil
yang berlawanan.

\[
\text{rasio beban} = \frac{858}{594} \approx 1{,}44,
\qquad
\text{selisih mutu} = 40 - 34 = 6 \text{ piksel}.
\]

Pada perbandingan ini, penghematan beban sebesar 31 persen harus dibayar dengan
kemunduran mutu sebesar 6 piksel, yaitu sekitar delapan belas persen. Kali ini
harga itu terlalu mahal.

Kedua kesimpulan itu tidak dapat diperoleh hanya dengan
melihat tabel, melainkan menuntut perhitungan rasio dan selisih. Cara berpikir
itulah yang membedakan penilaian (ranah C5) dari sekadar penerapan (ranah C3).

\begin{catatabox}
Perhatikan pula bahwa kolom \istilah{Jarak akhir} pada program juga tercatat,
tetapi nilainya tidak dipakai pada analisis di atas. Nilai itu hanya
menggambarkan keadaan pada bingkai terakhir, sehingga mudah berubah menurut
kebetulan posisi pemain pada saat itu. Untuk menilai mutu perilaku, dipakailah
jarak \textbf{rata-rata}, yang menggambarkan perilaku sepanjang permainan.
Memilih ukuran yang tepat merupakan bagian dari pekerjaan menilai.
\end{catatabox}

\subsection{Contoh Kedua: Merancang Skenario Uji}

\begin{contohbox}
Anda hendak menguji dua rancangan: rancangan A dengan selang hitung ulang
setiap bingkai, dan rancangan B dengan selang 0,3 detik. Keduanya diuji pada
ketiga skenario pada Tabel~\ref{tab:tiga-skenario}, dengan jalur pemain yang
sama untuk kedua rancangan.
\end{contohbox}

Rancangan pengujiannya dapat disusun sebagai tabel dua arah, sebagaimana
diperlihatkan pada Tabel~\ref{tab:rancangan-uji}.

\begin{table}[htbp]
\centering
\caption{Rancangan pengujian dua rancangan pada tiga skenario}
\label{tab:rancangan-uji}
{\setlength{\tabcolsep}{4pt}\footnotesize%
\begin{tabularx}{\textwidth}{@{}L{2.6cm} C{2.6cm} C{2.6cm} Y@{}}
\toprule
\headrow \textbf{Skenario} & \textbf{Rancangan A (setiap bingkai)} & \textbf{Rancangan B (0,3 detik)} & \textbf{Yang dibandingkan} \\
\midrule
Ruang terbuka & & & Jarak rata-rata dan banyaknya pencarian. \\
\addlinespace
Labirin & & & Jarak rata-rata, banyaknya pencarian, dan apakah jalurnya berubah-ubah tanpa sebab. \\
\addlinespace
Rintangan dinamis & & & Waktu yang diperlukan untuk memperbaiki jalur yang usang. \\
\bottomrule
\end{tabularx}}
\tabsource{Setiap sel pada tabel itu diisi dengan angka hasil pengukuran dan
kesan mutu perilakunya.}
\end{table}

Tiga hal perlu diperhatikan pada rancangan pengujian itu.

\begin{enumerate}
  \item \textbf{Jalur pemainnya sama pada kedua rancangan.} Inilah yang membuat
        hasilnya dapat dibandingkan. Apabila jalurnya berbeda, perbedaan
        hasilnya dapat berasal dari jalur pemain, bukan dari rancangan yang
        diuji.
  \item \textbf{Ketiga skenario diuji untuk kedua rancangan.} Perbaikan yang
        menguntungkan pada satu skenario dapat merugikan pada skenario lain.
        Rancangan yang baik adalah yang terbaik secara keseluruhan, bukan yang
        terbaik pada satu keadaan.
  \item \textbf{Setiap sel memuat angka dan kesan.} Angka untuk dibandingkan,
        kesan untuk menjelaskan. Laporan yang memuat keduanya jauh lebih
        berguna daripada laporan yang hanya memuat salah satunya.
\end{enumerate}
